\documentclass[11pt]{article}
\usepackage[T1]{fontenc}
\usepackage{textcomp}
\usepackage[margin=0.75in]{geometry}
\usepackage{amsmath,amssymb,amsthm}
\usepackage{array,booktabs}
\usepackage{hyperref}
\setlength{\parindent}{0pt}
\newtheorem{theorem}{Theorem}
\theoremstyle{definition}
\newtheorem{definition}{Definition}
\title{Flow Proof Artifact: Nat-order-plus-characterization}
\author{Generated by \texttt{flow doc proof}}
\date{Flow Proof Book}
\begin{document}
\maketitle
Less-or-equal is characterized by adding a natural on the right.\\[0.5em]
\textbf{Source.} peano — https://en.wikipedia.org/wiki/Total\_order\\[0.5em]
\bigskip
\noindent{\large \textbf{Derived fact 1.} \textit{a ≤ b if and only if there exists k with a + k = b (forward: addition witness)} \hfill the natural numbers \cdot order \cdot less-or-equal means a right summand exists}\par
\smallskip\noindent\textit{Coordinate: the natural numbers \cdot order \cdot less-or-equal means a right summand exists}\par
\smallskip\noindent\textit{Source: peano — Landau, *Foundations of Analysis*}\par
\smallskip\noindent\textit{Built on: adding on the right preserves order, for addition on the natural numbers}\par
\medskip
\noindent\textbf{Goal.} a ≤ b if and only if there exists k with a + k = b (forward: addition witness)\par
\noindent$\displaystyle \forall a \in \mathbb{N} \forall b \in \mathbb{N} \forall k \in \mathbb{N}\quad a \le b$\par
\medskip
\noindent
\begin{tabular}{@{} >{\raggedright\arraybackslash}p{0.06\textwidth} >{\raggedright\arraybackslash}p{0.47\textwidth} >{\raggedleft\arraybackslash}p{0.41\textwidth} @{}}
\textbf{\#} & \textbf{Proof} & \textbf{Mathematics} \\
\midrule
\textbf{1.} & We prove that less-or-equal means a right summand exists for order on the natural numbers. & \\[0.45em]
\textbf{2.} & We split into exhaustive cases — the claim must hold in each one. & \\[0.45em]
\textbf{3.} & Case 1 (see step 2): suppose a + k  equals  b. & \\[0.45em]
\textbf{4.} & We invoke the derived fact governing addition on the natural numbers: adding on the right preserves order, for addition on the natural numbers (instantiated for 0, a, b). & \\[0.45em]
\textbf{5.} & From step 3 and step 4, this implies a is at most b. Together with the other cases (step 3), the goal is discharged. Hence proven. & $\displaystyle a \le b$ \\[0.45em]
\bottomrule
\end{tabular}
\label{thm:Nat-order-less_or_equal_means_a_right_summand_exists}
\par\medskip\hrule
\end{document}
